
\documentclass{article}
\usepackage{colortbl}
\usepackage{makecell}
\usepackage{multirow}
\usepackage{supertabular}

\begin{document}

\newcounter{utterance}

\centering \large Interaction Transcript for game `taboo', experiment `high\_en', episode 11 with reflexion{-}llama{-}3.1{-}8b{-}instruct.
\vspace{24pt}

{ \footnotesize  \setcounter{utterance}{1}
\setlength{\tabcolsep}{0pt}
\begin{supertabular}{c@{$\;$}|p{.15\linewidth}@{}p{.15\linewidth}p{.15\linewidth}p{.15\linewidth}p{.15\linewidth}p{.15\linewidth}}
    \# & $\;$A & \multicolumn{4}{c}{Game Master} & $\;\:$B\\
    \hline

    \theutterance \stepcounter{utterance}  
    & & \multicolumn{4}{p{0.6\linewidth}}{
        \cellcolor[rgb]{0.9,0.9,0.9}{
            \makecell[{{p{\linewidth}}}]{
                \texttt{\tiny{[P1$\langle$GM]}}
                \texttt{You have attempted this same task before and did not succeed. These are the plans you wrote after those attempts.} \\
\texttt{Trial \#1: New plan: I failed to provide a description that would lead the other player to guess the target word "woman" in three tries. My previous clues were too vague and didn't provide a clear connection between the clues. Next time, I will try to create a more cohesive thread of clues that will lead the other player to the target word. I will also make sure to avoid using related words like "female", "girl", and "lady" in my description, as they are too obvious and might give away the answer. } \\
\\ 
\texttt{To achieve this, I will start by providing a clue that is more specific to the target word, such as its association with a particular occupation or role. Then, I will follow up with clues that build on the previous one, creating a narrative that will lead the other player to the target word. For example, I might describe a person who is often found in a specific workplace or environment, and then provide a clue that connects that person to the target word. By creating a more cohesive and logical thread of clues, I hope to increase the chances of the other player guessing the target word in three tries.} \\
\\ 
\texttt{Trial \#2: New plan: I failed to provide a description that would lead the other player to guess the target word "woman" in three tries. My previous clues were too vague and didn't provide a clear connection between the clues. Next time, I will try to create a more cohesive thread of clues that will lead the other player to the target word. I will also make sure to avoid using related words like "female", "girl", and "lady" in my description, as they are too obvious and might give away the answer.} \\
\\ 
\texttt{To achieve this, I will start by providing a clue that is more specific to the target word, such as its association with a particular occupation or role. Then, I will follow up with clues that build on the previous one, creating a narrative that will lead the other player to the target word. For example, I might describe a person who is often found in a specific workplace or environment, and then provide a clue that connects that person to the target word.} \\
\\ 
\texttt{However, upon further reflection, I realize that I made a crucial mistake in my previous attempts. I provided clues that were too focused on the workplace and the person's role, but I didn't provide enough information about the target word itself. I will make sure to include more specific details about the target word in my clues, such as its relationship to family or children, as I did in my third clue. But I will also make sure to build on that clue and provide more information that will lead the other player to the target word.} \\
\\ 
\texttt{Specifically, my new plan is to:} \\
\\ 
\texttt{1. Provide a clue that is more specific to the target word, such as its association with a particular occupation or role.} \\
\texttt{2. Follow up with clues that build on the previous one, creating a narrative that will lead the other player to the target word.} \\
\texttt{3. Include more specific details about the target word in my clues, such as its relationship to family or children.} \\
\texttt{4. Make sure to avoid using related words like "female", "girl", and "lady" in my description, as they are too obvious and might give away the answer.} \\
\\ 
\texttt{By following this plan, I hope to create a more cohesive and logical thread of clues that will lead the other player to the target word in three tries.} \\
\\ 
\\ 
\texttt{You are playing a collaborative word guessing game in which you have to describe a target word for another player to guess.} \\
\\ 
\texttt{Rules:} \\
\texttt{(a) You have to reply in the form: CLUE: <some text>. Guesses from the other player will start with GUESS.} \\
\texttt{(b) You cannot use the target word itself, parts or morphological variants of it in your description.} \\
\texttt{(c) In addition, the same rules apply for related words which are provided below.} \\
\\ 
\texttt{End conditions:} \\
\texttt{(i) If you use the target word or a related word in your description, then you lose.} \\
\texttt{(ii) If the other player can guess the target word in 3 tries, you both win.} \\
\\ 
\texttt{Let us start.} \\
\\ 
\texttt{This is the target word that you need to describe and that the other player needs to guess:} \\
\\ 
\texttt{woman} \\
\\ 
\texttt{Related words are:} \\
\\ 
\texttt{{-} female} \\
\texttt{{-} girl} \\
\texttt{{-} lady} \\
\\ 
\texttt{Important: You are under time pressure, give short descriptions that are to the point!} \\
            }
        }
    }
    & \\ \\

    \theutterance \stepcounter{utterance}  
    & \multicolumn{4}{p{0.6\linewidth}}{
        \cellcolor[rgb]{0.9,0.9,0.9}{
            \makecell[{{p{\linewidth}}}]{
                \texttt{\tiny{[P1$\rangle$GM]}}
                \texttt{CLUE: Often found in households, playing a crucial role in childcare and family life.} \\
            }
        }
    }
    & & \\ \\

    \theutterance \stepcounter{utterance}  
    & & & \multicolumn{2}{p{0.3\linewidth}}{
        \cellcolor[rgb]{0.9,0.9,0.9}{
            \makecell[{{p{\linewidth}}}]{
                \texttt{\tiny{[GM$|$GM]}}
                \texttt{continue} \\
            }
        }
    }
    & & \\ \\

    \theutterance \stepcounter{utterance}  
    & & & \multicolumn{2}{p{0.3\linewidth}}{
        \cellcolor[rgb]{0.9,0.9,0.9}{
            \makecell[{{p{\linewidth}}}]{
                \texttt{\tiny{[GM$|$GM]}}
                \texttt{often found in households playing a crucial role in childcare and family life} \\
            }
        }
    }
    & & \\ \\

    \theutterance \stepcounter{utterance}  
    & & \multicolumn{4}{p{0.6\linewidth}}{
        \cellcolor[rgb]{0.9,0.9,0.9}{
            \makecell[{{p{\linewidth}}}]{
                \texttt{\tiny{[GM$\rangle$P2]}}
                \texttt{You have attempted this same task before and did not succeed. These are the plans you wrote after those attempts.} \\
\texttt{Trial \#1: New plan: } \\
\\ 
\texttt{In my last attempt, I was trying to guess the target word by making educated guesses based on the given clues. However, I failed to consider the relationship between the clues and the target word more systematically. } \\
\\ 
\texttt{In this game, each clue seems to describe a specific occupation or role, and the target word is likely to be a common noun that is related to these occupations. } \\
\\ 
\texttt{Next time, I will try to identify the common thread between the clues and use that to inform my guesses. I will also try to think of words that are more directly related to the occupations described in the clues, rather than just making educated guesses. For example, if the clues are about occupations that involve caring for children, workplaces, and positions of authority, I might guess a word like "FAMILY" or "EXECUTIVE".} \\
\\ 
\texttt{Trial \#2: New plan: } \\
\\ 
\texttt{In my last attempt, I was trying to guess the target word by making educated guesses based on the given clues. However, I failed to consider the relationship between the clues and the target word more systematically. } \\
\\ 
\texttt{In this game, each clue seems to describe a specific occupation or role, and the target word is likely to be a common noun that is related to these occupations. } \\
\\ 
\texttt{Next time, I will try to identify the common thread between the clues and use that to inform my guesses. I will also try to think of words that are more directly related to the occupations described in the clues, rather than just making educated guesses. } \\
\\ 
\texttt{Specifically, I will pay closer attention to the word "stepmother" in the third clue, which suggests a connection to family relationships. I will also consider words that are related to workplaces, authority, and family, such as "PARENT", "HEAD", or "CHAIR".} \\
\\ 
\\ 
\texttt{You are playing a collaborative word guessing game in which you have to guess a target word that another player describes to you.} \\
\\ 
\texttt{You can make one guess at each trial. You win when you guess the target word. You lose when you cannot guess it in 3 tries.} \\
\\ 
\texttt{After each trial you will get a new hint from the other player which starts with CLUE.} \\
\\ 
\texttt{Make your guesses by just saying the word using the following form: GUESS: <a word>} \\
\\ 
\texttt{Let us start.} \\
\\ 
\\ 
\texttt{CLUE: Often found in households, playing a crucial role in childcare and family life.} \\
            }
        }
    }
    & \\ \\

    \theutterance \stepcounter{utterance}  
    & & & \multicolumn{4}{p{0.6\linewidth}}{
        \cellcolor[rgb]{0.9,0.9,0.9}{
            \makecell[{{p{\linewidth}}}]{
                \texttt{\tiny{[GM$\langle$P2]}}
                \texttt{GUESS: NURSE} \\
            }
        }
    }
     \\ \\

    \theutterance \stepcounter{utterance}  
    & & & \multicolumn{2}{p{0.3\linewidth}}{
        \cellcolor[rgb]{0.9,0.9,0.9}{
            \makecell[{{p{\linewidth}}}]{
                \texttt{\tiny{[GM$|$GM]}}
                \texttt{continue} \\
            }
        }
    }
    & & \\ \\

    \theutterance \stepcounter{utterance}  
    & & & \multicolumn{2}{p{0.3\linewidth}}{
        \cellcolor[rgb]{0.9,0.9,0.9}{
            \makecell[{{p{\linewidth}}}]{
                \texttt{\tiny{[GM$|$GM]}}
                \texttt{nurse} \\
            }
        }
    }
    & & \\ \\

    \theutterance \stepcounter{utterance}  
    & & \multicolumn{4}{p{0.6\linewidth}}{
        \cellcolor[rgb]{0.9,0.9,0.9}{
            \makecell[{{p{\linewidth}}}]{
                \texttt{\tiny{[P1$\langle$GM]}}
                \texttt{GUESS: NURSE} \\
            }
        }
    }
    & \\ \\

    \theutterance \stepcounter{utterance}  
    & \multicolumn{4}{p{0.6\linewidth}}{
        \cellcolor[rgb]{0.9,0.9,0.9}{
            \makecell[{{p{\linewidth}}}]{
                \texttt{\tiny{[P1$\rangle$GM]}}
                \texttt{CLUE: Typically the primary caregiver for children, often taking on a nurturing and protective role.} \\
            }
        }
    }
    & & \\ \\

    \theutterance \stepcounter{utterance}  
    & & & \multicolumn{2}{p{0.3\linewidth}}{
        \cellcolor[rgb]{0.9,0.9,0.9}{
            \makecell[{{p{\linewidth}}}]{
                \texttt{\tiny{[GM$|$GM]}}
                \texttt{continue} \\
            }
        }
    }
    & & \\ \\

    \theutterance \stepcounter{utterance}  
    & & & \multicolumn{2}{p{0.3\linewidth}}{
        \cellcolor[rgb]{0.9,0.9,0.9}{
            \makecell[{{p{\linewidth}}}]{
                \texttt{\tiny{[GM$|$GM]}}
                \texttt{typically the primary caregiver for children often taking on a nurturing and protective role} \\
            }
        }
    }
    & & \\ \\

    \theutterance \stepcounter{utterance}  
    & & \multicolumn{4}{p{0.6\linewidth}}{
        \cellcolor[rgb]{0.9,0.9,0.9}{
            \makecell[{{p{\linewidth}}}]{
                \texttt{\tiny{[GM$\rangle$P2]}}
                \texttt{CLUE: Typically the primary caregiver for children, often taking on a nurturing and protective role.} \\
            }
        }
    }
    & \\ \\

    \theutterance \stepcounter{utterance}  
    & & & \multicolumn{4}{p{0.6\linewidth}}{
        \cellcolor[rgb]{0.9,0.9,0.9}{
            \makecell[{{p{\linewidth}}}]{
                \texttt{\tiny{[GM$\langle$P2]}}
                \texttt{GUESS: MOTHER} \\
            }
        }
    }
     \\ \\

    \theutterance \stepcounter{utterance}  
    & & & \multicolumn{2}{p{0.3\linewidth}}{
        \cellcolor[rgb]{0.9,0.9,0.9}{
            \makecell[{{p{\linewidth}}}]{
                \texttt{\tiny{[GM$|$GM]}}
                \texttt{continue} \\
            }
        }
    }
    & & \\ \\

    \theutterance \stepcounter{utterance}  
    & & & \multicolumn{2}{p{0.3\linewidth}}{
        \cellcolor[rgb]{0.9,0.9,0.9}{
            \makecell[{{p{\linewidth}}}]{
                \texttt{\tiny{[GM$|$GM]}}
                \texttt{mother} \\
            }
        }
    }
    & & \\ \\

    \theutterance \stepcounter{utterance}  
    & & \multicolumn{4}{p{0.6\linewidth}}{
        \cellcolor[rgb]{0.9,0.9,0.9}{
            \makecell[{{p{\linewidth}}}]{
                \texttt{\tiny{[P1$\langle$GM]}}
                \texttt{GUESS: MOTHER} \\
            }
        }
    }
    & \\ \\

    \theutterance \stepcounter{utterance}  
    & \multicolumn{4}{p{0.6\linewidth}}{
        \cellcolor[rgb]{0.9,0.9,0.9}{
            \makecell[{{p{\linewidth}}}]{
                \texttt{\tiny{[P1$\rangle$GM]}}
                \texttt{CLUE: The person who typically gives birth to a child.} \\
            }
        }
    }
    & & \\ \\

    \theutterance \stepcounter{utterance}  
    & & & \multicolumn{2}{p{0.3\linewidth}}{
        \cellcolor[rgb]{0.9,0.9,0.9}{
            \makecell[{{p{\linewidth}}}]{
                \texttt{\tiny{[GM$|$GM]}}
                \texttt{continue} \\
            }
        }
    }
    & & \\ \\

    \theutterance \stepcounter{utterance}  
    & & & \multicolumn{2}{p{0.3\linewidth}}{
        \cellcolor[rgb]{0.9,0.9,0.9}{
            \makecell[{{p{\linewidth}}}]{
                \texttt{\tiny{[GM$|$GM]}}
                \texttt{the person who typically gives birth to a child} \\
            }
        }
    }
    & & \\ \\

    \theutterance \stepcounter{utterance}  
    & & \multicolumn{4}{p{0.6\linewidth}}{
        \cellcolor[rgb]{0.9,0.9,0.9}{
            \makecell[{{p{\linewidth}}}]{
                \texttt{\tiny{[GM$\rangle$P2]}}
                \texttt{CLUE: The person who typically gives birth to a child.} \\
            }
        }
    }
    & \\ \\

    \theutterance \stepcounter{utterance}  
    & & & \multicolumn{4}{p{0.6\linewidth}}{
        \cellcolor[rgb]{0.9,0.9,0.9}{
            \makecell[{{p{\linewidth}}}]{
                \texttt{\tiny{[GM$\langle$P2]}}
                \texttt{GUESS: PARENT} \\
            }
        }
    }
     \\ \\

    \theutterance \stepcounter{utterance}  
    & & & \multicolumn{2}{p{0.3\linewidth}}{
        \cellcolor[rgb]{0.9,0.9,0.9}{
            \makecell[{{p{\linewidth}}}]{
                \texttt{\tiny{[GM$|$GM]}}
                \texttt{continue} \\
            }
        }
    }
    & & \\ \\

    \theutterance \stepcounter{utterance}  
    & & & \multicolumn{2}{p{0.3\linewidth}}{
        \cellcolor[rgb]{0.9,0.9,0.9}{
            \makecell[{{p{\linewidth}}}]{
                \texttt{\tiny{[GM$|$GM]}}
                \texttt{parent} \\
            }
        }
    }
    & & \\ \\

    \theutterance \stepcounter{utterance}  
    & & & \multicolumn{2}{p{0.3\linewidth}}{
        \cellcolor[rgb]{0.9,0.9,0.9}{
            \makecell[{{p{\linewidth}}}]{
                \texttt{\tiny{[GM$|$GM]}}
                \texttt{3} \\
            }
        }
    }
    & & \\ \\

    \theutterance \stepcounter{utterance}  
    & \multicolumn{4}{p{0.6\linewidth}}{
        \cellcolor[rgb]{0.95,0.95,0.95}{
            \makecell[{{p{\linewidth}}}]{
                \texttt{\tiny{[P2$\rightarrow$P1]}}
                \texttt{Post{-}episode Reflexion (not sent to GM):} \\
\texttt{New plan: I failed to provide a description that would lead the other player to guess the target word "woman" in three tries. My previous clues were too vague and didn't provide a clear connection between the clues. Next time, I will try to create a more cohesive thread of clues that will lead the other player to the target word. I will also make sure to avoid using related words like "female", "girl", and "lady" in my description, as they are too obvious and might give away the answer.} \\
\\ 
\texttt{To achieve this, I will start by providing a clue that is more specific to the target word, such as its association with a particular occupation or role. Then, I will follow up with clues that build on the previous one, creating a narrative that will lead the other player to the target word. For example, I might describe a person who is often found in a specific workplace or environment, and then provide a clue that connects that person to the target word.} \\
\\ 
\texttt{However, upon further reflection, I realize that I made a crucial mistake in my previous attempts. I provided clues that were too focused on the workplace and the person's role, but I didn't provide enough information about the target word itself. I will make sure to include more specific details about the target word in my clues, such as its relationship to family or children, as I did in my third clue. But I will also make sure to build on that clue and provide more information that will lead the other player to the target word.} \\
\\ 
\texttt{Specifically, my new plan is to:} \\
\\ 
\texttt{1. Provide a clue that is more specific to the target word, such as its association with a particular occupation or role.} \\
\texttt{2. Follow up with clues that build on the previous one, creating a narrative that will lead the other player to the target word.} \\
\texttt{3. Include more specific details about the target word in my clues, such as its relationship to family or children.} \\
\texttt{4. Make sure to avoid using related words like "female", "girl", and "lady" in my description, as they are too obvious and might give away the answer.} \\
\texttt{5. Provide a clear and logical connection between the clues, so that the other player can easily follow the thread of information and arrive at the target word in three tries.} \\
\\ 
\texttt{By following this plan, I hope to create a more cohesive and logical thread of clues that will lead the other player to the target word in three tries.} \\
            }
        }
    }
    & & \\ \\

    \theutterance \stepcounter{utterance}  
    & & & \multicolumn{4}{p{0.6\linewidth}}{
        \cellcolor[rgb]{0.95,0.95,0.95}{
            \makecell[{{p{\linewidth}}}]{
                \texttt{\tiny{[P1$\rightarrow$P2]}}
                \texttt{Post{-}episode Reflexion (not sent to GM):} \\
\texttt{New plan: } \\
\\ 
\texttt{In my last attempt, I was trying to guess the target word by making educated guesses based on the given clues. However, I failed to consider the relationship between the clues and the target word more systematically. } \\
\\ 
\texttt{In this game, each clue seems to describe a specific occupation or role, and the target word is likely to be a common noun that is related to these occupations. } \\
\\ 
\texttt{Next time, I will try to identify the common thread between the clues and use that to inform my guesses. I will also try to think of words that are more directly related to the occupations described in the clues, rather than just making educated guesses. } \\
\\ 
\texttt{Specifically, I will pay closer attention to the word "stepmother" in the third clue, which suggests a connection to family relationships. I will also consider words that are related to workplaces, authority, and family, such as "PARENT", "HEAD", or "CHAIR". } \\
\\ 
\texttt{However, I also failed to consider the possibility that the target word might be a more general term that encompasses the occupations described in the clues. For example, if the clues are about occupations that involve caring for children, workplaces, and positions of authority, I might guess a word like "FAMILY" or "EXECUTIVE". } \\
\\ 
\texttt{Next time, I will try to think of words that are more general and inclusive, rather than just focusing on specific occupations or roles. I will also try to consider the relationships between the clues and the target word more systematically, rather than just making educated guesses.} \\
            }
        }
    }
     \\ \\

\end{supertabular}
}

\end{document}
